\section{One positive native factor for all original degrees and sections}

\begin{definition}\label{def:native-tc-one-proper-center-power-frame}
\lean{EnvelopeProperty.NativeTestConfiguration.TestConfiguration.gaussCompatibleCenterFrame,EnvelopeProperty.NativeTestConfiguration.TestConfiguration.gaussCompatibleTargetFrame}
\leanok
\uses{thm:native-tc-product-polynomial-gauss-source,thm:proper-model-norm-source-frame-existence,lem:scheme-line-actual-center-frame}
At the actual original product Gauss point, choose one genuine
unit frame of the original total line over its actual proper
center valuation ring. Restrict this one frame to the native
generic residue line and denote it by $e_M$. Its powers are the
genuine native pulled power frames of this same $e_M$.
\end{definition}
\begin{proof}
The product Gauss point has the actual polynomial Gauss base.
For $0<r\leq1$ this entire polynomial valuation is bounded by one,
as follows from the genuine polynomial Gauss formula and the
trivial coefficient valuation. Original total-space properness
therefore supplies its unique valuation-ring lift. The original
invertible line becomes free of rank one on this local valuation
ring, so a unit frame exists. Choose such a frame once and use
the actual native center-to-generic comparison to restrict it.
This constructs an original residue frame before any power or
section is chosen.
\end{proof}

\begin{theorem}\label{thm:native-tc-coherent-frame-all-section-norm}
\lean{EnvelopeProperty.NativeTestConfiguration.TestConfiguration.gaussCompatibleTargetFrame_power_norm_one,EnvelopeProperty.NativeTestConfiguration.TestConfiguration.gaussPowerSectionNorm_eq_compatibleFrame}
\leanok
\uses{def:native-tc-one-proper-center-power-frame,thm:fc-proper-center-power-frame-coherence,thm:proper-model-norm-framed-source,thm:native-tc-compatible-product-power-frames}
Every genuine power-frame inverse unit has true original
total-model norm one. For every $n\geq0$ and every original
section of $M^{\otimes n}$, its true total-model Gauss norm is
the native coefficient value of its actual entire product
transport in the model power residue frame constructed from
this one $e_M$.
\end{theorem}
\begin{proof}
The actual proper-center power frames restrict to exactly the
genuine generic power frames. Their inverse units have true
model norm one by the actual center-frame norm formula, in every
degree. The actual original product map sends the original
product Gauss point to the same total Gauss point. Apply genuine
framed pullback of proper-target norms through this actual map
and the literal original product power isomorphism. Its target
frame-unit factor is the preceding number one. The transported
whole section is exactly the original product section by its
definition, giving the asserted equality for all sections,
including zero, without properness of the Laurent source.
\end{proof}

\begin{definition}\label{def:native-tc-coherent-positive-norm-factor}
\lean{EnvelopeProperty.NativeTestConfiguration.TestConfiguration.gaussCompatibleNormFactor,EnvelopeProperty.NativeTestConfiguration.TestConfiguration.gaussCompatibleNormFactor_pos}
\leanok
\uses{def:native-tc-one-proper-center-power-frame,def:native-tc-compatible-product-residue-frames,thm:native-tc-frame-transition-unit-powers}
Choose an actual original frame $e_U$ of $L$ on a single original
open $U$ containing the actual proper center of $x$. Define $c$
to be the actual product residue valuation of the native unit
transition from its original reference residue frame to the
model residue frame induced from the preceding one $e_M$.
Then $c>0$, and this definition precedes every degree, section
and finite section family.
\end{definition}
\begin{proof}
The original product Gauss point lies in the inverse image of
$U$: the actual completed-point map factors through $U$, and
its genuine whole open product-point identity gives membership.
Thus the actual reference residue frame exists on the same
product residue line as the model residue frame. Their native
transition is a genuine unit, so its residue valuation is
strictly positive. Neither a section nor a degree occurs in
this construction.
\end{proof}

\begin{theorem}\label{thm:native-tc-coherent-all-power-section-norms}
\lean{EnvelopeProperty.NativeTestConfiguration.TestConfiguration.gaussPowerSectionNorm_eq_compatibleFactor,EnvelopeProperty.NativeTestConfiguration.TestConfiguration.exists_weightedLaurentPowerSectionNorm_powerFactor}
\leanok
\uses{thm:native-tc-coherent-frame-all-section-norm,def:native-tc-coherent-positive-norm-factor,thm:native-tc-frame-transition-unit-powers,cor:fc-native-laurent-support-frame-values,def:native-tc-weighted-laurent-section-norm}
For the actual factor $c$ just defined, every $n\geq0$ and every
original total-space section $s\in H^0(\mathcal X,M^{\otimes n})$
satisfy
\[
 \lVert s\rVert_{\mathcal X,r,x}=c^n
 \max_i\{\lVert s_i(x)\rVert_{L^{\otimes n}}r^i\},
\]
where $s_i$ are the actual entire original Laurent coefficients.
In particular genuine original invertibility of $L$ derives the
existence of one $c>0$ before all degrees and all original
sections. Negative exponents and zero sections are retained.
\end{theorem}
\begin{proof}
Compute the true total-model norm in the preceding compatible
model power frame. The genuine transition identity changes it
to the coefficient value in the original reference power frame,
multiplied by the native transition value. The actual transition
unit power theorem identifies this factor with $c^n$. The
reference power frame is the support frame obtained from the
same original $U$ and the genuine power frame of $e_U$.
Its native coefficient value is the actual Laurent support
maximum of the true original coefficient norms. This is the
displayed equality, simultaneously for every degree and every
section. Finally an original trivializing cover for the actual
invertible $L$ contains its actual proper center in one open;
choose that original open and frame once. The explicit positive
factor then proves the existential statement with the required
order of quantifiers.
\end{proof}

\begin{corollary}\label{cor:native-tc-coherent-all-power-family-norms}
\lean{EnvelopeProperty.NativeTestConfiguration.TestConfiguration.exists_weightedLaurentPowerNorms_powerFactor}
\leanok
\uses{thm:native-tc-coherent-all-power-section-norms,thm:native-tc-global-laurent-fixed-family}
The same one positive $c$ works for all degrees, all original
sections and all original finite section families. The maximum
total norm on a family equals $c^n$ times the maximum of true
weighted original coefficient norms on its actual finite
coefficient/exponent family, chosen before evaluation.
\end{corollary}
\begin{proof}
Choose $c$ from the preceding all-degree all-section identity.
For each degree and original finite family take its finite
maximum. Multiplication by the nonnegative $c^n$ preserves that
maximum. The maximum of the individual Laurent support maxima
is exactly the maximum over their original finite union, which
is the actual coefficient/exponent family. Empty families and
repeated values give the same equality, including zero.
\end{proof}
